\chapter*{Introduction}
\addcontentsline{toc}{chapter}{Introduction} % Ensures it appears in the Table of Contents
\pagestyle{plain}

The contamination of aquatic ecosystems by pharmaceutical compounds has become an urgent environmental challenge. Among the most promising remediation strategies is the use of specialized polymer sponges capable of absorbing these pollutants from wastewater. However, discovering the optimal polymer for a specific drug molecule is a complex and time--consuming process, traditionally driven by trial--and--error laboratory experiments. In recent years, computational methods and machine learning (ML) have emerged as powerful tools to accelerate the discovery of new materials \cite{AIPoweredMaterialScience}. This thesis explores a dual--track machine learning approach to address this challenge, focusing on both the generation of novel molecular structures and the prediction of polymer--molecule adsorption efficiency.

In line with the methodological priorities of this research, the thesis begins by introducing the fundamental machine learning architectures employed. First, we examine the Multi--Layer Perceptron (MLP), a robust feedforward neural network used for predictive tasks. Next, we present minGPT, a lightweight implementation of the Generative Pre--trained Transformer architecture, which serves as the backbone for generative modeling. The primary objective of deploying these models is to explore the complex chemical space of polymer--molecule interactions, using algorithmic pattern recognition to bridge the gap between chemical structures and physical properties.

To anchor these machine learning models in chemical reality, substantial effort was devoted to data collection and feature engineering. A comprehensive dataset was assembled by systematically extracting information from scientific literature on the use of polymers for the adsorption of molecules in wastewater. The key variables collected include the solution pH, the initial concentration of the pharmaceutical (mg/L), and the adsorption capacity (mg/g), defined as the amount of pollutant absorbed per unit mass of polymer.  Given the well--known scarcity of data describing these specific interaction triplets, the dataset was further augmented through interpolation techniques. In addition, to convert chemical information into a machine--readable representation, computational chemistry libraries — specifically RDKit, Polymetrix, and tblite — were employed to derive a rich set of features from the raw data.

The experimental phase of this work is structured around two main contributions. The first consists of a generative pipeline based on a custom wrapper of minGPT, designed to produce synthetic yet chemically valid polymer SMILES (Simplified Molecular Input Line Entry System) strings, thereby enabling a theoretical expansion of the candidate chemical space. The second contribution addresses the predictive task: employing the extracted features and a Multi--Layer Perceptron (MLP) to estimate the adsorption capacity of a given polymer--molecule pair. To assess the model’s ability to generalize under conditions of limited real--world data, extensive experiments were conducted using k--fold cross--validation in combination with systematic hyperparameter tuning. 

Ultimately, this thesis evaluates the current potential of machine learning approaches for predicting environmental remediation metrics, highlighting both the promise of molecular generation techniques and the persistent challenges associated with data scarcity in predictive chemical modeling.

The remainder of this thesis is organized as follows: Chapter 1 introduces the chemical context and selected features. Chapter 2 details the mathematical foundations of the MLP and minGPT architectures. Chapter 3 outlines the data scraping, interpolation, and feature extraction pipeline. Chapter 4 presents the generative modeling approach for P--SMILES. Chapter 5 discusses the predictive experiments, cross-validation results, and current limitations of the proposed approach. Finally, Chapter 6 provides concluding remarks and directions for future research.